
\chapter{Introducción} 

El presente documento detalla exhaustivamente el proceso de investigación y desarrollo 
llevado a cabo en el proyecto de fin de carrera de los estudiantes del programa de 
Ingeniería en Sistemas de Información de la Universidad Nacional de Costa Rica. 
Este proyecto se centra en la creación e implementación de soluciones de software 
innovadoras, diseñadas para equipar a los estudiantes con las habilidades necesarias 
para enfrentar los desafíos del entorno laboral actual en tecnologías de la información.

El objetivo primordial es proporcionar una experiencia práctica y relevante de integración en 
escenarios reales de trabajo, promoviendo así el desarrollo de competencias críticas como 
el análisis de problemas, diseño de sistemas, programación avanzada, gestión de proyectos y 
colaboración interdisciplinaria. A través de este enfoque, el proyecto no solo pretende 
reforzar el conocimiento teórico adquirido durante el curso académico, sino también fomentar 
la innovación y la creatividad en la solución de problemas complejos.

Asimismo, este documento expone los criterios metodológicos adoptados, incluyendo las fases 
de planificación, ejecución y evaluación del proyecto. Se describen las herramientas y tecnologías 
seleccionadas, justificando su uso en base a las tendencias actuales del mercado y su relevancia 
para los objetivos de aprendizaje del programa. Además, se discuten los resultados obtenidos, 
proporcionando análisis y conclusiones que destacan la contribución del proyecto al desarrollo 
profesional de los estudiantes.

En resumen, este documento sirve como un archivo comprensivo del trabajo realizado, también como un
 recurso formativo que guía a los estudiantes en la aplicación efectiva 
de sus habilidades técnicas en el mundo real, preparándolos para contribuir significativamente en 
sus futuros roles como profesionales de sistemas de información.
